% Part: computability
% Chapter: computability-theory
% Section: introduction

\documentclass[../../../include/open-logic-section]{subfiles}

\begin{document}

\olfileid[tr]{cmp}{thy}{int}
\olsection{Giriş}

Mantığın \emph{hesaplanabilirlik kuramı} diye bilinen dalı, fonksiyonlar ile
kümelerin hesaplanabilirliği ya da göreli hesaplanabilirliğiyle ilgili
sorunları inceler. Alanın eskiden \emph{özyineleme kuramı} diye bilinmesi
Kleene'nin etkisinin bir göstergesidir; günümüzde iki ad da yaygın biçimde
kullanılır.

Bir $f\colon\Nat\pto\Nat$ fonksiyonuna, belirli bir hesaplama modelinde
hesaplanabiliyorsa \emph{kısmi hesaplanabilir} diyelim. $f$ tümelse kısaca
$f$ \emph{hesaplanabilir} diyeceğiz. Karakteristik fonksiyonu $\Char{R}$
hesaplanabilir olan bir $R$ bağıntısına da hesaplanabilir denir. $f$ ile $g$
kısmi fonksiyonlarsa $f$'nin $x$'te tanımlı, yani $x$'in $f$'nin tanım
kümesinde olduğunu $f(x)\fdefined$ ile; $f$'nin $x$'te tanımlı olmadığını ise
$f(x)\fundefined$ ile göstereceğiz. $f(x)\simeq g(x)$, ya iki değerin de
tanımsız ya da ikisinin de tanımlı ve eşit olduğu anlamına gelecektir.

Belirli bir hesaplama modeline başvurmadan da hesaplanabilirlik kuramı
incelenebilir. Bunun için evrensel bir kısmi hesaplanabilir
$\fn{Un}(k,x)$ fonksiyonunun bulunduğu gösterilir. Bu, kısmi hesaplanabilir
fonksiyonları numaralandırmamızı sağlar. $k$. tek yerli kısmi hesaplanabilir
fonksiyonu göstermek üzere $\cfind{k}$ yazacak ve onu
$\cfind{k}(x)\simeq\fn{Un}(k,x)$ ile tanımlayacağız. Kleene bu amaçla
$\{k\}$ gösterimini kullanmıştı, fakat son zamanlarda bu gösterim daha az
kullanılmaktadır. Biraz daha genel olarak keyfî aritideki kısmi hesaplanabilir
fonksiyonları birbiçimli olarak numaralandırabiliriz; $k$. $n$-yerli kısmi
özyinelemeli fonksiyonu göstermek üzere $\cfind{k}[n]$ yazacağız.

$f(\vec x,y)$ tümel ya da kısmi bir fonksiyonsa
$\umin{y}{f(\vec x,y)}$, $f(\vec x,0),\dots,f(\vec x,y-1)$ değerlerinin
tümünün tanımlı olması koşuluyla $f(\vec x,y)=0$ olacak en küçük $y$'yi
döndüren $\vec x$ fonksiyonudur; böyle bir $y$ yoksa
$\umin{y}{f(\vec x,y)}$ tanımsızdır. $R(\vec x,y)$ bir bağıntıysa
$\umin{y}{R(\vec x,y)}$, $R(\vec x,y)$ doğru olacak en küçük $y$, başka bir
deyişle $1\tsub\Char{R}(\vec x,y)=0$ olacak en küçük $y$ olarak tanımlanır.

Bir fonksiyonun hesaplanabilir olduğunu göstermek için iki yoldan biri
izlenebilir:
\begin{enumerate}
\item Titiz yol: Bir Turing makinesini ya da kısmi özyinelemeli fonksiyonu
  açıkça betimleyip amaçlanan fonksiyonu hesapladığını göstermek.
\item Biçimsel olmayan yol: Fonksiyonu hesaplayan bir algoritmayı betimleyip
  Church tezine başvurmak.
\end{enumerate}
İkisi arasında keskin bir sınır yoktur. Bir algoritmanın ayrıntılı betimi,
uygun Turing makinesinin ya da kısmi özyinelemeli tanımlar dizisinin ilkece
nasıl tasarlanabileceğini yeterince açık kılmalıdır. Bütünüyle biçimsel
tanımların aydınlatıcı olması beklenmez. Bu iki yaklaşım arasında uygun bir
denge bulmaya, kısacası kesin tanımların elde edilebileceğini gösteren hem
sezgisel hem titiz ispatlar vermeye çalışacağız.

\end{document}

